\begin{textsample}{NEG Dim 4 – persona\_grok – Score -1.00 – t210\_grok.txt}  \label{ex:f4_neg_026}
The climate crises of 2021, marked by devastating heat waves and floods, have been described as a"code red for humanity."These events have disproportionately impacted developing nations, though their effects are now felt worldwide.
To address this urgency, ending corporate greed is essential.
This involves redirecting resources from the super-rich away from destructive practices and toward climate solutions, such as by \textbf{taxing} excessive wealth.
Key strategies include banning advertisements for fossil fuels, supporting Indigenous resistance efforts, and accelerating the transition to renewable energy sources.
A group of 183,000 ultra-wealthy individuals holds outsized influence, and \textbf{taxing} them can enable the systemic changes needed to combat the climate crisis.
This year marks 50 years since the publication of"The Limits to Growth,"yet society continues on an unsustainable path, exacerbated by rising \textbf{inequality}.
Greenpeace is partnering with the Fight \textbf{Inequality} Alliance to prioritize the well-being of the planet over private profit.
The \textbf{inequalities} exposed by COVID-19 offer a preview of the disparities that climate change will bring.

% matched lemmas: inequality, tax
\end{textsample}
